problem:  
Let \( (M = G/H, g) \) be a compact, simply connected, **naturally reductive** Riemannian manifold with reductive decomposition \(\mathfrak{g} = \mathfrak{h} \oplus \mathfrak{m}\) and canonical connection \(\tilde{\nabla}\) with torsion \(T(X,Y) = -[X,Y]_{\mathfrak m}\).  

For each connection \( \nabla'\in\{\nabla,\tilde{\nabla}\}\), define  
\[
\mathfrak{r}^{\nabla'} := \mathrm{Span}\{ R^{\nabla'}(X,Y) : X,Y\in \mathfrak{m}\} \subset \mathfrak{so}(\mathfrak{m}),
\]
the **curvature operator span**, that is, the smallest linear subspace of \(\mathfrak{so}(\mathfrak{m})\) generated by all curvature endomorphisms of \(\nabla'\).  
Write also \(\mathfrak{h}_c := \mathrm{Lie}\langle \mathfrak{r}^{\tilde{\nabla}}\rangle = \mathfrak{hol}(\tilde{\nabla})\) and \(\mathfrak{h}_\ell := \mathrm{Lie}\langle \mathfrak{r}^{\nabla}\rangle = \mathfrak{hol}(\nabla)\) for the canonical and Levi-Civita holonomy algebras, respectively.

**Open Problem (Curvature–Span Rigidity Conjecture):**  
Assume that the **linear spans** of the corresponding curvature operators coincide,
\[
\mathfrak{r}^{\nabla} = \mathfrak{r}^{\tilde{\nabla}} \subset \mathfrak{so}(\mathfrak{m}).
\]
Must it follow that \((M,g)\) is **normal homogeneous**, i.e. that the naturally reductive metric arises from a bi-invariant metric on a larger transitive Lie group?

Equivalently: *Does equality of curvature operator spans for the Levi-Civita and canonical connections characterize normal homogeneity among naturally reductive metrics?*

---

Why is it a "good" problem:  

1. **Mathematically precise and nontrivial:**  
   The curvature operator span \(\mathfrak{r}^{\nabla'}\) is a rigorously defined subspace of \(\mathfrak{so}(\mathfrak{m})\) encoding all curvature directions generated by the connection. Equality of these spans is a well-posed algebraic condition distinct from both tensor equality and holonomy equality. It imposes a meaningful linear constraint rather than a vacuous one.

2. **Bridges algebraic curvature and geometric structure:**  
   The condition \(\mathfrak{r}^{\nabla} = \mathfrak{r}^{\tilde{\nabla}}\) equates the curvature-generation patterns of the torsion-free and skew-torsion connections. Resolving whether this forces normal homogeneity would clarify the extent to which canonical torsion actually contributes new curvature directions, connecting algebraic curvature operators with large-scale group symmetries.

3. **Conceptual depth via minimal assumptions:**  
   Among all possible rigidity hypotheses, equating curvature spans is a remarkably weak algebraic statement—yet it may suffice to determine a highly structured geometric outcome (normal homogeneity). This tension between a minimal assumption and a potentially strong conclusion makes it profoundly probing.

4. **Interplay with holonomy theory:**  
   The problem sits at an intersection of **homogeneous curvature algebra**, **connection holonomy**, and **representation theory of isotropy groups**. It naturally extends questions of Ambrose–Singer type (holonomy generation by curvature) to the specific comparative context of two canonical connections on naturally reductive spaces.

5. **Testable in concrete families:**  
   Known naturally reductive examples—such as Berger spheres, flag manifolds, nearly Kähler manifolds, and certain Stiefel manifolds—permit explicit computation of both \(\mathfrak{r}^{\nabla}\) and \(\mathfrak{r}^{\tilde{\nabla}}\). Checking whether equality occurs provides a concrete testing ground for the conjecture and possible classification of any nontrivial cases.

6. **Potential geometric payoff:**  
   A positive resolution would provide an algebraic–curvature characterization of normal homogeneous metrics within the broader naturally reductive class. A counterexample would expose families of geometries where canonical torsion is “curvature-invisible,” revealing unexpected subtleties between curvature generation and symmetry.

7. **Simplicity with depth:**  
   The formulation  
   \[
   \mathfrak{r}^{\nabla} = \mathfrak{r}^{\tilde{\nabla}} \quad \Rightarrow \quad (M,g)\ \text{is normal homogeneous?}
   \]
   is succinct, self-contained, and accessible to experts in homogeneous geometry, yet resolving it demands a synthesis of Lie-algebraic, holonomic, and geometric insight. This clarity of expression coupled with latent depth embodies the ideal of the “narrow entrance, profound depth” principle of a good problem.